% =============================================================================
% Drop-in replacement for results.tex lines 69-88.
% Generated 2026-05-10 in response to Comments 1, 4, 5, 6, 7.
%
% Companion artifacts produced by this update:
%   - scripts/python/21_realized_vs_anticipated.py
%       -> figures/realized_vs_anticipated.png
%   - scripts/python/22_belief_specs.py
%       -> tables/reg_belief_specs.tex   (label: tab:reg_belief_specs)
%
% The existing tab:del_decision_determinants (tables/reg_delegation.tex)
% is left UNCHANGED. The new belief table is a separate artifact with its
% own label.
% =============================================================================

\paragraph{Punishment beliefs as a candidate channel.} A natural reading of Result~1 is that Player~A anticipates stronger punishment for a delegated decision than for a self-made decision and delegates less in response. This belief pattern would be consistent with parts of prior work suggesting that observers view delegating moral decisions to algorithms relatively critically \citep[e.g.,][]{gogoll_rage_2018, niszczota_robo-investors_2020, jauernig_people_2022}. Alternatively, the same belief pattern would be consistent with a simple shirking motive, in which Player~A expects greater punishment for shirking the effort to make a decision for Player~B but, in the absence of punishment, can shirk without monetary consequences. We have already seen that Player~B punishment is outcome-based, rather than delegation-based (which is not surprising given that performance is controlled for), but a shirking motive --- which would predict extra punishment for delegation --- could at least rationalise the positive punishment observed in the good outcome--delegated scenario. In either case, delegation behaviour should be consistent with observed beliefs about punishment.

\paragraph{Punishment beliefs.} \label{sec:beliefs} After making the delegation decision but before learning the outcome, we elicited Player~A's beliefs about Player~B's punishment in each of the four scenarios.

Several caveats may temper the interpretation of the belief measure. The elicitation was unincentivised and took place after the delegation decision, so reports may reflect hedging across scenarios or post-hoc rationalisation rather than the beliefs that actually drove the choice. As a corollary, any within-subject correlation between beliefs and delegation cannot identify direction: we cannot distinguish whether subjects delegated \textit{because} they held these specific beliefs about punishment, or hold these beliefs \textit{because} they delegated. Hypotheticality compounds the noise: only two of the four \textit{(delegation, outcome)} scenarios remained payoff-relevant after the delegation decision in the \textit{Punishment} condition, and all four were purely hypothetical in the \textit{No-Punishment} condition. We therefore restrict the belief--delegation analysis to the \textit{Punishment} condition. The motivation is structural rather than statistical: the punishment-anticipation channel can only operate where punishment is possible, so beliefs about a counterfactual that cannot materialise do not enter Player~A's decision rule in the \textit{No-Punishment} arm. The exercise that follows should accordingly be read as suggestive evidence on whether the channel is alive among those exposed to it, not as a mediation analysis of the treatment-level delegation gap (Result~1).

Two patterns emerge. First, Player~As substantially \textit{overestimate} the level of punishment they will face: average expected punishment exceeds average realised punishment in every scenario (Figure~\ref{fig:realized_vs_anticipated}, green vs.\ magenta bars). The same pattern is visible in the \textit{No-Punishment} condition, where Player~A's hypothetical beliefs (blue bars) closely track the Punishment-condition beliefs --- consistent with the noise sources flagged above, but reassuring as a basic calibration check on the elicitation. Second, however, Player~As correctly anticipate the \textit{relative} pattern of punishment: they expect no significant difference in punishment between delegated and non-delegated decisions, conditional on outcome --- a pattern that closely tracks the actual punishment behaviour of Player~B documented in Result~2. \textit{It is this relative-pattern finding --- on punishment beliefs alone --- that provides no support for a punishment-anticipation mechanism for Result~1.} If Player~A had reduced delegation in the \textit{Punishment} condition because she anticipated being punished more harshly for delegation than for self-made decisions, we would expect to see this asymmetry in the elicited beliefs. We do not.

\begin{figure}[!htbp]
    \centering
    \caption{Realised vs.\ anticipated punishment, by scenario}
    \includegraphics[width=0.85\linewidth]{figures/realized_vs_anticipated.png}
    \label{fig:realized_vs_anticipated}
    \begin{minipage}{0.9\linewidth}
    \footnotesize
    \textit{Note:} Average punishment in pounds across the four \textit{(delegation, outcome)} scenarios, attention-check passers only. Green bars: Player~A's anticipated punishment in the \textit{Punishment} condition ($n=61$). Magenta bars: Player~B's realised punishment in the \textit{Punishment} condition ($n=67$). Blue bars: Player~A's hypothetical punishment beliefs in the \textit{No-Punishment} condition ($n=70$), shown for calibration only --- these beliefs concern a counterfactual that cannot materialise and do not enter Player~A's delegation decision in this arm. Whiskers denote $\pm 1$ standard error.
    \end{minipage}
\end{figure}

Even though the belief measure was not incentivised, and even though at the time of the belief elicitation two of the four scenarios could not materialise anymore, punishment beliefs carry within-subject information. Table~\ref{tab:reg_belief_specs} reports a Logit regression of the delegation decision on Player~A's belief difference (\textit{no delegation}~$-$~\textit{delegation}) in expected punishment, restricted throughout to the \textit{Punishment} condition.\footnote{The two underlying belief measures compare expected punishment of a self-made and delegated decision for each of the two possible outcomes for Player~B, so that higher values correspond to higher expected punishment for non-delegated decisions relative to delegated decisions.} The first three columns use the simple average of the two outcome-conditional belief differences as the single regressor; this maps cleanly onto a punishment-anticipation channel under expected-utility reasoning, in which only the expected punishment penalty for switching from delegation to self-decision matters. Column~(1) reports the unconditional specification on the full \textit{Punishment} sample; Column~(2) adds task performance and socio-demographic controls; Column~(3) restricts to subjects who passed the attention check on the belief-elicitation screen, as preregistered. Column~(4) replaces the simple average with the more theoretically appropriate weighted average, using Player~A's elicited probability of producing the high payoff as the weight ($\Pr(\text{good})=$\textit{wa\_confidence}$/10$). Column~(5) instead enters the two outcome-conditional belief differences separately, as in earlier drafts.

The pattern is consistent with the punishment-anticipation channel \textit{conditional on attention}, but unconditionally weak. In Columns~(1)--(2), the belief coefficient is positive but small and far from significant ($p>0.39$): in the unrestricted \textit{Punishment} sample, expected differential punishment does not predict delegation. Restricting to attention-check passers (Column~(3)) sharpens the relationship substantially: the simple-average belief coefficient is~$2.71$ ($p=0.035$), implying that subjects who expect a one-pound larger punishment penalty for not delegating are markedly more likely to delegate. The weighted-difference specification (Column~(4)) yields a similar conclusion with a slightly larger coefficient ($3.33$, $p=0.030$) and the highest pseudo-$R^2$ in the table, supporting the expected-utility interpretation. In Column~(5), the good-outcome belief difference carries the regression: its coefficient of~$1.58$ ($p=0.033$) is in line with the single-regressor specifications, while the bad-outcome belief difference is positive but statistically noisier ($0.94$, $p=0.21$).

\begin{table}[!htbp]
\centering
\scriptsize
\caption{Delegation and punishment beliefs}
\label{tab:reg_belief_specs}
\begin{tabular}{@{\extracolsep{5pt}}lcc|ccc}
\\[-1.8ex]\hline
\hline \\[-1.8ex]
& \multicolumn{5}{c}{Dependent variable: \textit{Delegation}}\\[0.1cm]
& \multicolumn{2}{c}{\textit{Full Punishment}} & \multicolumn{3}{c}{\textit{Punishment, passers}} \\
\\[-1.8ex] & (1) & (2) & (3) & (4) & (5) \\
\hline \\[-1.8ex]
 Punishment beliefs: simple average ($\overline{\Delta}$) & $0.553$ & $0.629$ & $2.708^{**}$ &  &  \\
  & $(0.656)$ & $(0.744)$ & $(1.286)$ &  &  \\[0.1cm]
 Punishment beliefs: weighted by $\Pr(\text{good})$ &  &  &  & $3.329^{**}$ &  \\
  &  &  &  & $(1.532)$ &  \\[0.1cm]
 Punishment beliefs: bad outcome (no del.~$-$~del.) &  &  &  &  & $0.937$ \\
  &  &  &  &  & $(0.751)$ \\[0.1cm]
 Punishment beliefs: good outcome (no del.~$-$~del.) &  &  &  &  & $1.579^{**}$ \\
  &  &  &  &  & $(0.741)$ \\[0.1cm]
 Task performance &  & $-0.337^{*}$ & $-0.576^{**}$ & $-0.544^{**}$ & $-0.609^{***}$ \\
  &  & $(0.181)$ & $(0.231)$ & $(0.245)$ & $(0.236)$ \\[0.1cm]
 Age &  & $0.037^{*}$ & $0.063^{**}$ & $0.066^{**}$ & $0.064^{**}$ \\
  &  & $(0.020)$ & $(0.029)$ & $(0.030)$ & $(0.030)$ \\[0.1cm]
 Female &  & $0.213$ & $1.522^{**}$ & $1.630^{**}$ & $1.334^{*}$ \\
  &  & $(0.592)$ & $(0.754)$ & $(0.699)$ & $(0.793)$ \\[0.1cm]
 Socio-economic status &  & $-0.079$ & $-0.012$ & $-0.095$ & $-0.003$ \\
  &  & $(0.185)$ & $(0.226)$ & $(0.233)$ & $(0.219)$ \\[0.1cm]
 Went to uni &  & $0.108$ & $-0.446$ & $-0.257$ & $-0.433$ \\
  &  & $(0.609)$ & $(0.712)$ & $(0.719)$ & $(0.707)$ \\[0.1cm]
 Technology score &  & $-0.104$ & $0.085$ & $0.129$ & $0.059$ \\
  &  & $(0.230)$ & $(0.290)$ & $(0.295)$ & $(0.283)$ \\[0.1cm]
 Leadership position &  & $0.478$ & $-0.485$ & $-0.477$ & $-0.387$ \\
  &  & $(0.510)$ & $(0.720)$ & $(0.731)$ & $(0.729)$ \\[0.1cm]
 Constant & $-0.316$ & $-0.536$ & $-2.023$ & $-2.153$ & $-1.842$ \\
  & $(0.228)$ & $(1.333)$ & $(2.080)$ & $(2.052)$ & $(2.103)$ \\[0.1cm]
\hline \\[-1.8ex]
 Controls & No & Yes & Yes & Yes & Yes \\
 Sample & \textit{Punishment} & \textit{Punishment} & Passers & Passers & Passers \\
 N & $80$ & $78$ & $60$ & $60$ & $60$ \\
 Pseudo $R^2$ & $0.006$ & $0.097$ & $0.241$ & $0.258$ & $0.246$ \\
\hline
\hline \\[-1.8ex]
\multicolumn{6}{p{0.85\textwidth}}{\footnotesize \textit{Note:} Logit estimates with robust (HC1) standard errors in parentheses. Significance: $^{*}\,p<0.10$; $^{**}\,p<0.05$; $^{***}\,p<0.01$ (two-sided). Sample restricted to the \textit{Punishment} condition throughout, since beliefs about a counterfactual that cannot materialise do not enter Player~A's decision rule in the \textit{No-Punishment} arm. Belief differences are within-subject \textit{(no delegation)}~$-$~\textit{(delegation)} in expected punishment, so positive coefficients indicate that subjects who expect to be punished more harshly for not delegating (relative to delegating) are more likely to delegate. Column~(1) uses the simple average of the two outcome-conditional belief differences without controls; Column~(2) adds task performance and socio-demographic controls; Column~(3) restricts to subjects who passed the attention check on the belief-elicitation screen, as preregistered; Column~(4) replaces the simple average with a weighted average using Player~A's elicited probability of producing the high payoff ($\Pr(\text{good})=$ \textit{wa\_confidence}$/10$); Column~(5) replaces the single belief regressor with the two outcome-conditional belief differences entered separately.}
\end{tabular}
\end{table}


\paragraph*{Interpretation of the two-difference decomposition.}
The five-column table treats the simple- and weighted-average specifications as the primary structural mappings of the channel and the two-difference decomposition (Column~(5)) as a descriptive complement. Two interpretive points are worth flagging. First, the two outcome-conditional differences are nearly orthogonal in this sample (Pearson correlation $-0.10$ in the passers sample), so the two coefficients in Column~(5) are essentially additive partial slopes rather than competing for the same variance. Second, under expected-utility reasoning only the weighted sum of the two differences should enter the decision rule; entering both separately therefore answers the partial question ``conditional on the bad-outcome asymmetry, how does delegation respond to the good-outcome asymmetry?'' rather than the structural question the channel poses. We therefore lead with the simple-average and weighted-average specifications and treat the decomposition as a robustness check that confirms the action sits in the good-outcome cell.

The within-subject signal is asymmetric across cells (Figure~\ref{fig:belief_distributions}): the single significant difference between delegators and non-delegators is in the \textit{self decision + good outcome} cell, where delegators expect punishment of $\pounds 0.82$ while non-delegators expect only $\pounds 0.32$ ($t$-test $p=0.006$, Mann-Whitney $p=0.019$); for the other three cells the gap is small and not significant. Because beliefs were elicited after the delegation decision and the elicitation was unincentivised, we cannot rule out that subjects rationalised their choice by holding consistent post-hoc beliefs; the within-subject correlation is consistent with the standard interpretation that anticipated punishment shapes delegation, but does not identify the causal direction.

\begin{figure}[!htbp]
    \centering
    \caption{Player A beliefs about Player B punishment, by delegation choice}
    \includegraphics[width=0.85\linewidth]{figures/belief_distributions.png}
    \label{fig:belief_distributions}
    \begin{minipage}{0.9\linewidth}
    \footnotesize
    \textit{Note:} Distribution of Player~A's expected punishment from Player~B in each of the four (delegation, outcome) cells, separately for Player~As who actually delegated (blue) and those who did not (orange). \textit{Punishment} condition, attention-check passers ($n_{\text{deleg}}=24$, $n_{\text{nodeleg}}=37$). The only cell with a statistically significant difference is \textit{self decision + good outcome} (top right): delegators expected $\pounds 0.82$ vs non-delegators $\pounds 0.32$ ($t$-test $p=0.006$).
    \end{minipage}
\end{figure}
